% =====================================================================
% Section 2 — Background
% =====================================================================

\section{Background}\label{sec:background}

\subsection{Valve-train failure modes at high rotational speed}
\label{sec:bg-failure-modes}

At elevated rotational speeds, the valve train of an SI engine moves
from a quasi-static loading regime, in which the valve simply tracks the
cam profile, into an inertia-dominated regime in which the spring force
is the only mechanism that keeps the follower in contact with the cam
\citep{heywood1988,taylor1985}. The classical failure modes that emerge
in this regime are the following.

\paragraph{Valve float}
When the spring force is no longer sufficient to overcome the inertial
force of the moving mass on the deceleration ramp of the cam, the
follower separates from the cam. The valve then moves freely under the
combined effect of the spring force and the residual inertia, and
re-contact with the cam occurs in an uncontrolled manner. Float reduces
the effective valve-opening profile and introduces large impacts at
re-contact.

\paragraph{Valve bounce}
Bounce is the impact-driven oscillation of the valve at seating. At high
rpm the closing velocity grows with the cam acceleration, and the
combination of spring stiffness and seat compliance can result in
multiple short re-openings before the valve finally seats. Bounce
degrades sealing, accelerates seat wear, and disturbs the gas-exchange
process.

\paragraph{Coil bind}
Coil bind is the condition in which adjacent coils of the valve spring
are in metal-to-metal contact. When this happens, the spring stops
behaving as a linear element: stiffness rises sharply and the dynamic
load is transferred directly through the coil stack rather than through
the spring deformation. Even an instantaneous approach to coil bind can
produce a load spike sufficient to damage the spring, the retainer, or
the cam lobe.

\paragraph{Spring fatigue}
The valve spring operates under fluctuating shear loading at the cam
frequency, oscillating between a non-zero seat-preload stress and the
full-lift stress; the alternating component is therefore superimposed
on a positive mean stress and is more demanding than a fully-reversed
loading of the same amplitude. At 10\,000~rpm a single spring
undergoes of the order of $10^{8}$ cycles in 200~hours of operation.
Fatigue resistance therefore governs the long-term life, and any
local stress concentration (surface defect, shot-peening discontinuity,
corrosion) can dominate the failure process \citep{shigley2020}.

\paragraph{Resonance with cam harmonics}
The cam profile contains harmonics at integer multiples of the camshaft
rotational frequency. If a harmonic coincides with the natural frequency
of the spring, the spring enters resonance, the dynamic deflection grows
beyond the design displacement, and coil bind and surge become likely
\citep{heywood1988}.

These five modes are not independent: float can trigger an impact that
exceeds the residual coil-bind clearance, and resonance can amplify both.
The BowTie methodology, introduced next, provides a unified language to
describe how each mode contributes to a single top event.

% ----------------------------------------------------------------------
\subsection{The BowTie methodology}\label{sec:bg-bowtie}

The BowTie methodology organises a risk analysis around a single top
event. Threats (causes) are arranged on the left side of the bow,
connected to the top event through preventive barriers; consequences are
arranged on the right side, connected to the top event through
mitigative barriers. The diagram resembles a bowtie, hence the name. The
methodology has its origins in process-safety analysis and was
consolidated as a general risk-management framework by the review of
\citet{deruijter2016bowtiereview}; it has since been applied to
mechanical-reliability problems where the chain of causation deserves
explicit visualisation \citep{aust2019bowtieaircraft}.

Compared with FMEA, BowTie offers three advantages relevant to the
present problem:

\begin{enumerate}
\item \emph{Causal chain}: BowTie distinguishes threats, preventive
barriers, consequences, and mitigative barriers as separate elements;
FMEA collapses this into a single severity--occurrence--detectability
triplet.
\item \emph{Barrier-centric view}: each preventive barrier can be
associated with a measurable indicator and a pass criterion, supporting
the quantitative coupling proposed in this work.
\item \emph{Communication}: the diagrammatic representation is
accessible to designers, operators, and reliability engineers
simultaneously, which is valuable in the context of a master's research
that aims to bridge mechanical design and reliability analysis.
\end{enumerate}

The classical BowTie is qualitative. To compensate for that limitation,
several quantitative extensions have been proposed; in particular,
\citet{khakzad2013bowtiebayesian} map the BowTie into a Bayesian
network so that the diagram becomes dynamic and probabilistic, and the
review of \citet{deruijter2016bowtiereview} catalogues the alternative
extensions ranging from layer-of-protection analysis (LOPA) to
fault-tree-based barrier reliability. The approach
adopted here is more pragmatic: each preventive barrier of the diagram
is mapped to a deterministic indicator computed from the analytical
sizing of the component, with a literature-anchored pass
criterion. This avoids the need for probabilistic data that is not
available for the bespoke high-performance cylinder head under study.

% ----------------------------------------------------------------------
\subsection{Related work}\label{sec:bg-related}

Several authors have addressed valve-train and engine-component
reliability from complementary perspectives. The work of
\citet{velezgodino2019overheadvalvetrain} on overhead valve trains in
urban buses is a representative case study published in
\textit{Engineering Failure Analysis}: the authors use field data and
metallurgical inspection to attribute the recurrent failure of an
in-service valve train to a combination of operational misuse and
inadequate inspection intervals, but they stop short of organising the
findings around a barrier-based diagram. \citet{soffritti2018wornvalvetrain}
report a four-cylinder diesel engine in which excessive wear at the
rocker-arm/pushrod and rocker-arm/valve interfaces appeared after only
1\,000~hours of operation; their root-cause analysis points to valve
misalignment with respect to the seat insert as the dominant trigger.
Both studies confirm that valve-train failures, once initiated,
propagate quickly into other components and motivate the need for a
preventive analysis framework.

Two further studies in \textit{Engineering Failure Analysis} address
elements of the valve-train system that are central to the present
work. \citet{xing2021valvespringEAC} investigate environment-assisted
cracking in high-strength valve springs and show that the fatigue limit
adopted at the design stage must be discounted to account for surface
condition and lubricant chemistry; this finding directly supports the
fatigue factor of safety adopted in
Section~\ref{sec:method-bowtie-quant}.
\citet{zhao2023camshaftV6} report lubrication failure in the camshaft
bearings of a V6 diesel engine, identifying an asperity-contact regime
that violates the design assumption of full hydrodynamic lubrication
and reduces the effective service life of the bearing pair; their
result reinforces the importance of the lubrication-related preventive
barriers shown on the left side of the BowTie diagram of
Section~\ref{sec:method-bowtie-qual}.
\citet{ko2022springfatigueflaws} extend the spring-fatigue discussion
beyond \textit{Engineering Failure Analysis}, showing in
\textit{Scientific Reports} that surface flaws as small as a few
micrometres on the oil-tempered wire of an automotive valve spring can
reduce the fatigue life by an order of magnitude; this result motivates
the conservative factor of safety adopted in
Table~\ref{tbl:quant-criteria}.

Table~\ref{tbl:related-cases} consolidates the four
\textit{Engineering Failure Analysis} case studies above and maps
their reported root causes to the preventive (P) and mitigative (M)
barriers and the deterministic indicators of the BowTie procedure
constructed in Section~\ref{sec:methodology}. The mapping is direct
for two of the four cases (Xing on environment-assisted cracking,
caught by the Goodman fatigue indicator; and the operational/
maintenance signature of V\'elez~Godi\~no, caught by the mitigative
side of the BowTie); is partial for the assembly-driven misalignment
of Soffritti, which the qualitative bow flags but which lies outside
the deterministic-indicator set adopted here; and is explicitly
\emph{outside} the present scope for the camshaft-bearing
lubrication failure of Zhao, which would require an additional
preventive barrier on the lubrication regime. The mixed mapping is
itself an outcome of the comparison: it identifies which classes of
failure the procedure already covers and which classes call for an
extension of the barrier set, and it is examined in detail in
Section~\ref{sec:res-litvalidation}.

\begin{table*}[width=\textwidth,pos=t]
\caption{Mapping of four documented valve-train and engine-component failure case studies from \textit{Engineering Failure Analysis} to the preventive (P) and mitigative (M) barriers and the deterministic indicators of the BowTie procedure proposed in this work. The notation P1--P5, M1--M5 refers to the barriers introduced in Section~\ref{sec:method-bowtie-qual}, and the indicators are those of Table~\ref{tbl:quant-criteria}. Reported root causes are paraphrased from the cited papers; no values are reproduced.}\label{tbl:related-cases}
\renewcommand{\arraystretch}{1.15}
\begin{tabular}{@{}p{0.19\textwidth}p{0.26\textwidth}p{0.26\textwidth}p{0.20\textwidth}@{}}
\toprule
Case study (component / engine) & Reported root cause & BowTie barrier triggered & Indicator that would flag it \\
\midrule
\citet{velezgodino2019overheadvalvetrain}; overhead valve train, urban bus & Operational misuse and inadequate inspection intervals & M3 (operating envelope at design) and M5 (predictive maintenance via vibration) & Mitigative side; outside the deterministic indicator set \\
\addlinespace
\citet{soffritti2018wornvalvetrain}; rocker-arm / pushrod interfaces, four-cylinder diesel ($\sim$1\,000~h) & Valve misalignment relative to the seat insert & Threat: out-of-tolerance assembly $\rightarrow$ candidate geometric/quality barrier P6 (extension) & Outside the present indicator set; flagged qualitatively by the bow chain \\
\addlinespace
\citet{xing2021valvespringEAC}; high-strength valve springs (premature fatigue) & Environment-assisted cracking due to surface condition and lubricant chemistry & P4 (high-fatigue-resistance spring material) & Goodman factor of safety $n_{f}$ with discounted shear endurance limit $S_{se}$ \\
\addlinespace
\citet{zhao2023camshaftV6}; camshaft bearings, V6 diesel & Asperity-contact regime violating the full-hydrodynamic-lubrication assumption & Threat: lubrication degradation; \emph{outside the present preventive set} & Outside scope; motivates a lubrication-regime barrier in procedure extensions \\
\bottomrule
\end{tabular}
\end{table*}

The methodological literature on the BowTie diagram itself has
matured in parallel. The review by \citet{deruijter2016bowtiereview}
remains the canonical reference for the qualitative formulation and
catalogues the variations adopted in the chemical, aerospace, and
nuclear industries. \citet{khakzad2013bowtiebayesian} propose a mapping
of the BowTie into a Bayesian network so that the static qualitative
diagram becomes dynamic and probabilistic; their construction is
adopted in part by the BowTie tool used in the offshore industry.
\citet{aust2019bowtieaircraft} apply the methodology to aircraft engine
borescope inspection and show that, even in the absence of
probabilistic data, the BowTie diagram is an effective communication
artefact between maintenance planners, inspectors, and reliability
engineers.

It is also relevant to acknowledge a parallel branch of the literature
that extends the conventional FMEA in different directions: fuzzy-FMEA
formulations and analytic-hierarchy-process-based weighting (AHP-FMEA)
have been used in mechanical reliability to mitigate the loss of
resolution of the multiplicative Risk Priority Number (RPN). The
present work is complementary to those extensions: it does not aim
to refine the FMEA score but to replace the single-score logic with
a barrier-by-barrier deterministic check, and is therefore
positioned closer to the BowTie tradition than to fuzzy/AHP-FMEA.

To the best of the authors' knowledge, no published study applies the
BowTie methodology to the valve train of a high-rpm spark-ignition
cylinder head, and no study couples the BowTie diagram to the
analytical sizing of the intake valve and the valve spring as a
reproducible barrier-evaluation procedure. The present work fills both
gaps.
